\section[B. 左灯右行]{\myboxii{      \BF{\TT{B}}} 左灯右行}

\textit{运行时间限制}   1 second

\textit{内存使用限制} 1 Gigabyte

\textit{读入文件} \qquad 标准输入

\textit{输出文件} \qquad 标准输出


\begin{quote}
	\textit{我爱你有种}
\end{quote}

\BF{这是一道交互题。}

\subsection{题目描述}

小 T 面对 T 城的地铁图犯了难 ...

具体而言，这座城市在 E 河两岸，北岸有 $n$ 个地铁站点 $N_1$, $N_2$, $\ldots$, $N_n$, 南岸有 $n$ 个地铁站点 $S_1$, $S_2$, $\ldots$, $S_n$.

T 城总共有 $n - 1$ 条地铁线路，每条线路都是环线，其中第 $i$ 条环线连接的地铁站是 $N_i - N_{i + 1} - S_{i + 1} - S_i - N_i$, 但这些地铁环线都是单向运行的，可能是 $N_i \leftarrow N_{i + 1} \leftarrow S_{i + 1} \leftarrow S_i \leftarrow N_i$ 也可能是 $N_i \rightarrow N_{i + 1} \rightarrow S_{i + 1} \rightarrow S_i \rightarrow N_i$.


每条环线都有一个价目表 $u_i$, $d_i$, $l_i$, $r_i$, 分别表示这 4 个相邻站 ($N_i - N_{i + 1}$, $S_i - S_{i + 1}$, $N_i - S_i$ 和 $N_{i + 1} - S_{i + 1}$) 之间坐一站地铁移动的价格，注意坐一站地铁的方向由地铁的运行方向而定，可能是 $N_i \rightarrow N_{i + 1}$ 也可能是 $N_{i} \leftarrow N_{i + 1}$.

遗憾的是，时空黑洞吞噬了价目表上这些环线的运行方向和价格，小 T 准备求助大 K 的帮助来得知这些地铁的运行方向。

大 K 知道地铁的运行方向，不过，小 T 可以为每条环线的四段相邻站边重新任意指定一组价格；大 K 知道真实的运行方向，并会根据这些指定价格回答小 T 的询问。

小 T 可以询问他至多 $3\,500$ 个问题：

\begin{itemize}
	\item 地铁站 $x$(可以是 $N_i$ 也可以是 $S_i$, 下略) 到地铁站 $y$ 坐地铁，在小 T 指定的地铁价目表下，至少需要支付多少地铁费用？
   \item 若不同环线连接同一对相邻站，则这些边同时存在并且都可以使用。一次旅行的费用为所经过边权之和，站点间换乘不额外收费。
\end{itemize}


在询问完这些问题后，小 T 需要和大 K 确认每条地铁的运行方向。

请帮助小 T 解决这个问题。

请注意，地铁线路的方向是预先确定的，它们不会在交互的过程中更改，换言之交互器\BF{不是}自适应的。

\subsection{交互协议}

本题每个测试点只有一组测试数据。

第一行输入一个整数 $n$ ($2 \le n \le 10^5$), 表示参数，其中地铁站共 $2n$ 座，环线共 $n - 1$ 条。

你首先需要输出 $n - 1$ 行，其中第 $i$ 行输出 $4$ 个整数 $u_i$, $d_i$, $l_i$, $r_i$, 分别表示第 $i$ 条环线上 $N_i - N_{i + 1}$, $S_i - S_{i + 1}$, $N_i - S_i$ 和 $N_{i + 1} - S_{i + 1}$ 这四个相邻站之间的移动价格，你需要保证这四个整数在 $[0, 10^9]$ 之间。

随后你将会直接开始交互，每一次交互你可以以下面的格式向交互器询问问题：

\begin{itemize}
   \item 输出一行，首先输出一个字符 \TT{?} 表示你询问一次，随后输出 $4$ 个元素 $S_x$, $id_x$, $S_y$, $id_y$, 其中你需要保证 $S_x, S_y \in \{\mathtt{N}, \mathtt{S}\}$, $id_x, id_y \in \{m \in \mathbb{N} : 1 \le m \le n\}$, 表示小 T 需要询问大 K 位于 $S_x$ 岸的第 $id_x$ 个地铁站 $(S_x)_{id_x}$ 到位于 $S_y$ 岸的第 $id_y$ 个地铁站 $(S_y)_{id_y}$ 之间至少需要支付多少地铁费用。
   \item 若你的输入合法并且没有超出询问限制，交互器会给出一行一个整数 $d$ 供输入表示这两个站点之间最少需要的地铁费用。
   \item 若你询问的问题超出了对应的询问次数限制，交互器会输出一个整数 $-1$ 并结束你的交互进程，此时你会获得\textit{答案错误}的评测结果。
\end{itemize}

若你已经得到地铁方向的话，请按照以下的格式输出答案：

\begin{itemize}
   \item 输出一行，首先输出一个字符 \TT{!}, 随后输出一个长度为 $n - 1$ 的字符串 $S$, 其中第 $i$ 个字符 $S_i \in \{\mathtt{I}, \mathtt{O}\}$, 若 $S_i = \mathtt{I}$ 则表示第 $i$ 条环线的方向为 $N_i \leftarrow N_{i + 1} \leftarrow S_{i + 1} \leftarrow S_i \leftarrow N_i$, 若 $S_i = \mathtt{O}$ 则表示第 $i$ 条环线的方向为 $N_i \rightarrow N_{i + 1} \rightarrow S_{i + 1} \rightarrow S_i \rightarrow N_i$.
\end{itemize}

在完成输出答案后，你应当安全退出。

每次输出都需要输出换行并且刷新缓冲区，否则你可能会得到意料之外的结果。

为了刷新缓冲区，你可以：

\begin{itemize}
   \item 对于 C 或者 C++, 使用 \TT{fflush(stdout)} 或 \TT{cout.flush()}。
   \item 对于 Java 或 Kotlin, 使用 \TT{System.out.flush()}。
   \item 对于 Python, 使用 \TT{stdout.flush()}。
\end{itemize}

\subsection{样例 1}

\begin{minipage}{\linewidth}
	\begin{minipage}[t]{0.49\linewidth}
		\begin{lstlisting}[title=\texttt{standard input}]
3
10
12
\end{lstlisting}
	\end{minipage}
	\begin{minipage}[t]{0.49\linewidth}
		\begin{lstlisting}[title=\texttt{standard output}]
10 3 1 2
10 10 1 1
? N 1 N 2
? N 2 N 3
! OI
\end{lstlisting}
	\end{minipage}
\end{minipage}


\subsection{样例 1 解释}

\begin{figure}[htbp]
   \centering
\begin{tikzpicture}[
  >=Stealth,
  station/.style={circle,draw,thick,minimum size=10mm,font=\large},
  bank/.style={thick},
  solid/.style={thick,->},
  rdashed/.style={thick,dashed,->},
  edgecost/.style={font=\Large, inner sep=0pt}
]

% --- coordinates / stations ---
\node[station] (N1) at (0, 2.2) {$N_1$};
\node[station] (N2) at (4, 2.2) {$N_2$};
\node[station] (N3) at (8, 2.2) {$N_3$};

\node[station] (S1) at (0,-2.2) {$S_1$};
\node[station] (S2) at (4,-2.2) {$S_2$};
\node[station] (S3) at (8,-2.2) {$S_3$};

% --- banks (text NOT on the line) ---
\draw[bank] (-1.5, 0.9) -- (9.5, 0.9);
\node[font=\Large] at (9.4, 1.25) {North bank};

\draw[bank] (-1.5,-0.9) -- (9.5,-0.9);
\node[font=\Large] at (9.4,-0.55) {South bank};

% --- Ring 1 = O : N1->N2->S2->S1->N1 ---
\draw[solid] (N1) -- (N2) node[midway,above=8pt,edgecost] {10};
\draw[solid] (N2) -- (S2) node[midway,left=10pt,edgecost] {2};


\draw[solid] (S2) -- (S1) node[midway,below=8pt,edgecost] {3};
\draw[solid] (S1) -- (N1) node[midway,left=10pt,edgecost] {1};

% --- Ring 2 = I : N3->N2, S3->N3, S2->S3, N2->S2 ---
\draw[rdashed] (N3) -- (N2) node[midway,above=8pt,edgecost] {10};
\draw[rdashed] (S3) -- (N3) node[midway,right=10pt,edgecost] {2};
\draw[rdashed] (S2) -- (S3) node[midway,below=8pt,edgecost] {10};

% --- the requested "double edge" on N2->S2: one solid + one dashed ---
% dashed copy slightly shifted to the right, so both are visible
\draw[rdashed] ([xshift=3pt]N2.south) -- ([xshift=3pt]S2.north)
  node[midway,right=14pt,edgecost] {1};

% --- ring labels ---
\node[font=\Large] at (2.0, 0) {Ring 1};
\node[font=\Large] at (6.0, 0) {Ring 2};

\end{tikzpicture}


\end{figure}
上面的图展示了地铁的运行方向和小 T 的价格标注。

请注意样例中的交互过程仅供参考，实际交互过程并不唯一，且这一示例的交互过程不一定是可行的或是最优的，本题的样例不会在附加文件中出现。





